% BHU PYQ -- Geography: Economic Geography (2025-26 NEP)
\documentclass[11pt,a4paper]{article}
\usepackage[a4paper,top=2.5cm,bottom=2.5cm,left=2.8cm,right=2.8cm,headheight=14pt]{geometry}
\usepackage{amsmath,amssymb}
\usepackage{fontspec}
\setmainfont{Kohinoor Devanagari}
\usepackage{microtype,fancyhdr,enumitem,hyperref,lastpage,parskip}

\hypersetup{colorlinks=true,urlcolor=black,linkcolor=black,pdftitle={GGRMJ31: Economic Geography -- BHU NEP 2025-26}}

\pagestyle{fancy}\fancyhf{}
\renewcommand{\headrulewidth}{0.4pt}\renewcommand{\footrulewidth}{0.4pt}
\fancyhead[L]{\small \textit{Banaras Hindu University}}
\fancyhead[R]{\small \textit{GGRMJ31: Economic Geography (2025-26)}}
\fancyfoot[C]{\small Page \thepage\ of \pageref{LastPage}}

\newlist{parts}{enumerate}{2}
\setlist[parts,1]{label=(\alph*),leftmargin=*,itemsep=4pt,topsep=3pt}
\setlist[parts,2]{label=(\roman*),leftmargin=*,itemsep=2pt,topsep=2pt}

\newcommand{\pts}[1]{\hfill \textit{[#1]}}

\begin{document}

\begin{center}
  {\large\bfseries BANARAS HINDU UNIVERSITY}\\[2pt]
  {\normalsize B. A. / B. Sc. (Hons.) Semester III Examination, 2025-26 (NEP)}\\[6pt]
  \rule{0.8\linewidth}{0.5pt}\\[6pt]
  {\large\bfseries GEOGRAPHY}\\[4pt]
  {\large\bfseries Paper Code: GGRMJ31 : Economic Geography}\\[4pt]
  {\normalsize Time: 2:30 Hours \hfill Full Marks: 70}\\[2pt]
  {\small REF NO. CECONF/N/2025-26/086}
\end{center}
\rule{\linewidth}{0.4pt}
\smallskip

\noindent \textit{Instructions: Answer five questions in all including Question No.\ 1 which is compulsory. Answer each part of Question No.\ 1 in approximately 200 words and remaining questions in approximately 1000 words. Illustrate your answers with suitable maps and diagrams. Write your Roll No.\ at the top immediately on receipt of this question paper.}

\smallskip
\noindent \textit{प्रश्न संख्या 1 (अनिवार्य) सहित कुल पाँच प्रश्नों के उत्तर दीजिए। प्रश्न 1 के प्रत्येक भाग का उत्तर लगभग 200 शब्दों में एवं शेष प्रश्नों का उत्तर लगभग 1000 शब्दों में दीजिए। अपने उत्तरों को उपयुक्त मानचित्रों तथा आरेखों द्वारा स्पष्ट कीजिए।}

\bigskip

\begin{enumerate}[label=\textbf{\arabic*.},leftmargin=*,itemsep=14pt]

  \item Write short notes on the following:\pts{4 $\times$ 3.5 = 14}
  
  निम्नलिखित पर संक्षिप्त टिप्पणियाँ लिखिए:
  \begin{parts}
    \item Fishing as primary activity / प्राथमिक गतिविधि के रूप में मत्स्य पालन
    \item Special Economic Zones (SEZ) / विशेष आर्थिक क्षेत्र
    \item Scope of Economic Geography / आर्थिक भूगोल का विषय-क्षेत्र
    \item Types of industries / उद्योगों के प्रकार
  \end{parts}

  \item Discuss in detail, the Von Thünen theory of agricultural location.\pts{14}
  
  वॉन थ्यूनेन के कृषि अवस्थिति सिद्धान्त की विस्तार से चर्चा कीजिए।

  \item Discuss the forest resource classification of the World.\pts{14}
  
  विश्व के वन संसाधनों के वर्गीकरण की विवेचना कीजिए।

  \item Present a classification of agricultural regions of the World developed by D. Whittlesey and explain the main characteristics of any two of them.\pts{14}
  
  डी० ह्विटलेसी द्वारा विकसित विश्व के कृषि प्रदेशों का वर्गीकरण प्रस्तुत कीजिए तथा उनमें से किन्हीं दो की मुख्य विशेषताओं की व्याख्या कीजिए।

  \item Describe the salient features and importance of Trans-Continental Railways of Russia (Trans-Siberian Railway).\pts{14}
  
  रूस के ट्रांस-कॉन्टिनेंटल रेलवे (ट्रांस-साइबेरियन रेलमार्ग) की मुख्य विशेषताओं एवं महत्त्व का वर्णन कीजिए।

  \item Explain Weber's theory of industrial location.\pts{14}
  
  वेबर के औद्योगिक अवस्थिति सिद्धान्त की व्याख्या कीजिए।

  \item Discuss the spatial pattern of global trade and major international trade routes.\pts{14}
  
  वैश्विक व्यापार के स्थानिक प्रतिरूप तथा प्रमुख अन्तरराष्ट्रीय व्यापार मार्गों की चर्चा कीजिए।

\end{enumerate}

\end{document}
